\documentclass[../../main.tex]{subfiles}
\begin{document}
\section{向量丛}\label{section:ch10-向量丛}

设 $M$ 为一拓扑空间。$M$ 上一个秩为 $k$ 的（实）向量丛是一个拓扑空间 $E$ 连同一个满连续映射 $\pi\colon E\to M$，满足下列条件：
\begin{enumerate}[(i)]
\item 对每个 $p\in M$，$p$ 上的纤维 $E_p=\pi^{-1}(p)$ 被赋予一个 $k$ 维实向量空间的结构；
\item 对每个 $p\in M$，存在 $M$ 中 $p$ 的一个邻域 $U$ 以及一个同胚 $\Phi\colon \pi^{-1}(U)\to U\times\mathbb{R}^k$（称为 $E$ 在 $U$ 上的一个\textbf{局部平凡化}），满足下列条件（\cref{figure:ch10-局部平凡化}）：
\begin{itemize}
\item $\pi_U\circ\Phi=\pi$（其中 $\pi_U\colon U\times\mathbb{R}^k\to U$ 为投影）；
\item 对每个 $q\in U$，$\Phi$ 在 $E_q$ 上的限制是从 $E_q$ 到 $\{q\}\times\mathbb{R}^k\cong\mathbb{R}^k$ 的向量空间同构。
\end{itemize}
\end{enumerate}

\begin{figure}[htbp]
\centering
\includegraphics[width=0.8\textwidth]{fig-10-1.png}
\caption{向量丛的一个局部平凡化}\label{figure:ch10-局部平凡化}
\end{figure}

若 $M$ 与 $E$ 是带边或不带边的光滑流形，$\pi$ 是光滑映射，且局部平凡化可取为微分同胚，则 $E$ 称为\textbf{光滑向量丛}。此时，我们把任意到其像上为微分同胚的局部平凡化称为\textbf{光滑局部平凡化}。

秩为 $1$ 的向量丛常称为（实）线丛。复向量丛可类似地定义，只需在定义中把“实向量空间”换成“复向量空间”、把 $\mathbb{R}^k$ 换成 $\mathbb{C}^k$。本书无需讨论复向量丛，因此此后我们所说的向量丛均指实向量丛，不再另行说明。

空间 $E$ 称为丛的\textbf{全空间}，$M$ 称为其\textbf{底}，$\pi$ 称为其\textbf{投影}。根据我们希望强调的方面，有时会从记号中略去某些成分，写作“$E$ 是 $M$ 上的向量丛”，或“$E\to M$ 是向量丛”，或“$\pi\colon E\to M$ 是向量丛”。

\begin{exercise}\label{exercise:ch10-01}
设 $E$ 为 $M$ 上的光滑向量丛。证明投影映射 $\pi\colon E\to M$ 是一个满的光滑淹没。
\end{exercise}

若 $E$ 在整个 $M$ 上都存在局部平凡化（称为 $E$ 的\textbf{整体平凡化}），则 $E$ 称为\textbf{平凡丛}。此时 $E$ 本身与乘积空间 $M\times\mathbb{R}^k$ 同胚。若 $E\to M$ 是一个光滑丛且容许一个光滑整体平凡化，则称 $E$ 是\textbf{光滑平凡的}。此时 $E$ 与 $M\times\mathbb{R}^k$ 微分同胚，而不仅是同胚。为简便起见，当我们说一个光滑丛是平凡的，总是指光滑平凡，而不仅是拓扑意义下的平凡。

\begin{figure}[htbp]
\centering
\includegraphics[width=0.85\textwidth]{fig-10-2.png}
\caption{Möbius 丛的一部分}\label{figure:ch10-mobius丛}
\end{figure}

\begin{example}[乘积丛]\label{example:ch10-乘积丛}
任意空间 $M$ 上秩为 $k$ 的向量丛的一个特别简单的例子是乘积空间 $E=M\times\mathbb{R}^k$，以 $\pi=\pi_1\colon M\times\mathbb{R}^k\to M$ 为其投影。任何这样的丛称为\textbf{乘积丛}，是平凡的（以恒等映射为整体平凡化）。若 $M$ 是带边或不带边的光滑流形，则 $M\times\mathbb{R}^k$ 是光滑平凡的。
\end{example}

尽管存在许多非平凡的向量丛，但唯一易于可视化的是下面这一个。

\begin{example}[Möbius 丛]\label{example:ch10-mobius丛}
在 $\mathbb{R}^2$ 上定义一个等价关系：$(x,y)\sim(x',y')$ 当且仅当对某个 $n\in\mathbb{Z}$ 有 $(x',y')=(x+n,(-1)^n y)$。令 $E=\mathbb{R}^2/\sim$ 表示商空间，$q\colon\mathbb{R}^2\to E$ 为商映射。

为了可视化 $E$，令 $S$ 表示条带 $[0,1]\times\mathbb{R}\subseteq\mathbb{R}^2$。$q$ 在 $S$ 上的限制是满的且为闭映射，因而是商映射。$q|_S$ 所做的唯一非平凡等同发生在两条边界线上，因此我们可以把 $E$ 设想为由 $S$ 得到的空间：把右边沿扭转半圈使其上下颠倒，然后将其粘贴到左边沿上（\cref{figure:ch10-mobius丛}）。对任意 $r>0$，商映射 $q$ 把矩形 $[0,1]\times[-r,r]$ 映成的像，是一个光滑的紧带边流形，称为\textbf{Möbius 带}；你可以用纸条粘贴末端并扭转半圈来制作它的纸模型。

考虑下面的交换图：
\begin{align*}
\begin{array}{ccc}
\mathbb{R}^2 & \xrightarrow{q} & E\\
\pi_1\downarrow && \downarrow\pi\\
\mathbb{R} & \xrightarrow{\varepsilon} & S^1,
\end{array}
\end{align*}
其中 $\pi_1$ 是到第一分量的投影，$\varepsilon\colon\mathbb{R}\to S^1$ 是光滑覆叠映射 $\varepsilon(x)=e^{2\pi\mathrm{i}x}$。由于 $\varepsilon\circ\pi_1$ 在每个等价类上为常数，它下降为一个连续映射 $\pi\colon E\to S^1$。一个直接（虽然繁琐）的验证表明，$E$ 具有唯一的光滑流形结构使得 $q$ 是光滑覆叠映射，且 $\pi\colon E\to S^1$ 是 $S^1$ 上的光滑实线丛，称为\textbf{Möbius 丛}。（若 $U\subseteq S^1$ 是被 $\varepsilon$ 均匀覆叠的开子集，而 $\widetilde U\subseteq\mathbb{R}$ 是 $\varepsilon^{-1}(U)$ 的一个连通分支，则 $q$ 限制为从 $\widetilde U\times\mathbb{R}$ 到 $\pi^{-1}(U)$ 的同胚。利用这一点，可以构造从 $\pi^{-1}(U)$ 到 $U\times\mathbb{R}$ 的同胚，作为 $E$ 的局部平凡化。这些局部平凡化可解释为定义 $E$ 上光滑结构的坐标图。习题~\ref{exercise:ch10-05} 要求你补出细节。习题~15-9 则提示了一种更强有力的方法。）
\end{example}

向量丛最重要的例子是切丛。

\begin{proposition}[切丛作为向量丛]\label{proposition:ch10-切丛作为向量丛}
设 $M$ 为带边或不带边的光滑 $n$ 流形，$TM$ 为其切丛。连同其标准投影映射、每条纤维上的自然向量空间结构，以及在\cref{proposition:ch03-18}~中构造的拓扑与光滑结构，$TM$ 是 $M$ 上秩为 $n$ 的光滑向量丛。
\end{proposition}

\begin{proof}
任取 $M$ 的一个光滑坐标图 $(U,\varphi)$，其坐标函数为 $(x^i)$，定义映射 $\Phi\colon\pi^{-1}(U)\to U\times\mathbb{R}^n$ 为
\begin{align}
\Phi\left(\left.v^i\frac{\partial}{\partial x^i}\right|_p\right)=(p,(v^1,\dots,v^n)).\label{align:ch10-平凡化坐标}
\end{align}
它在纤维上是线性的，且满足 $\pi_1\circ\Phi=\pi$。复合映射
\begin{align*}
\pi^{-1}(U)\xrightarrow{\Phi}U\times\mathbb{R}^n\xrightarrow{\varphi\times\mathrm{Id}_{\mathbb{R}^n}}\varphi(U)\times\mathbb{R}^n
\end{align*}
等于在\cref{proposition:ch03-18}~中构造的坐标映射 $\widetilde\varphi$。由于 $\widetilde\varphi$ 与 $\varphi\times\mathrm{Id}_{\mathbb{R}^n}$ 都是微分同胚，$\Phi$ 亦然。因此 $\Phi$ 满足光滑局部平凡化的全部条件。
\end{proof}

任何非平凡的丛当然都需要不止一个局部平凡化。下面的引理表明，两个光滑局部平凡化在重叠处的复合具有简单的形式。

\begin{lemma}\label{lemma:ch10-转移函数}
设 $\pi\colon E\to M$ 为 $M$ 上秩为 $k$ 的光滑向量丛。设 $\Phi\colon\pi^{-1}(U)\to U\times\mathbb{R}^k$ 与 $\Psi\colon\pi^{-1}(V)\to V\times\mathbb{R}^k$ 是 $E$ 的两个光滑局部平凡化，且 $U\cap V\neq\varnothing$。存在光滑映射 $\gamma\colon U\cap V\to GL(k,\mathbb{R})$，使得复合 $\Phi\circ\Psi^{-1}\colon(U\cap V)\times\mathbb{R}^k\to(U\cap V)\times\mathbb{R}^k$ 具有形式
\begin{align}
\Phi\circ\Psi^{-1}(p,v)=(p,\gamma(p)v),\label{align:ch10-转移函数形式}
\end{align}
其中 $\gamma(p)v$ 表示 $k\times k$ 矩阵 $\gamma(p)$ 对向量 $v\in\mathbb{R}^k$ 的通常作用。
\end{lemma}

\begin{proof}
下面的图交换：
\begin{align*}
\begin{array}{ccccc}
(U\cap V)\times\mathbb{R}^k & \xleftarrow{\Psi} & \pi^{-1}(U\cap V) & \xrightarrow{\Phi} & (U\cap V)\times\mathbb{R}^k\\
& \pi_1\searrow & & \swarrow\pi_1\\
&& U\cap V
\end{array}
\end{align*}
（上图上方的映射应理解为 $\Psi$ 与 $\Phi$ 在 $\pi^{-1}(U\cap V)$ 上的限制。）由此可得 $\pi_1\circ(\Phi\circ\Psi^{-1})=\pi_1$，这意味着
\begin{align*}
\Phi\circ\Psi^{-1}(p,v)=(p,\eta(p,v))
\end{align*}
对某个光滑映射 $\eta\colon(U\cap V)\times\mathbb{R}^k\to\mathbb{R}^k$ 成立。此外，对每个固定的 $p\in U\cap V$，映射 $v\mapsto\eta(p,v)$ 是从 $\mathbb{R}^k$ 到其自身的可逆线性映射，因此存在一个非奇异的 $k\times k$ 矩阵 $\gamma(p)$ 使得 $\eta(p,v)=\gamma(p)v$。剩下只需证明映射 $\gamma\colon U\cap V\to GL(k,\mathbb{R})$ 是光滑的。这留作习题~\ref{exercise:ch10-04}。
\end{proof}

本引理所描述的光滑映射 $\gamma\colon U\cap V\to GL(k,\mathbb{R})$ 称为局部平凡化 $\Phi$ 与 $\Psi$ 之间的\textbf{转移函数}。（这是光滑流形理论中少数几个虽值域不是 $\mathbb{R}$ 或 $\mathbb{R}^k$ 却传统上仍用“函数”一词的情形之一。）例如，若 $M$ 是光滑流形，而 $\Phi$ 与 $\Psi$ 是与两个不同光滑坐标图相联系的 $TM$ 的局部平凡化，则坐标转移映射的 Jacobian 矩阵就是它们之间的转移函数。

与切丛一样，向量丛通常最容易通过给出一个向量空间的族来描述，即在底流形的每一点给出一个向量空间。为了把这样的集合变成光滑向量丛，我们首先要在所有向量空间的不交并上构造一个流形拓扑与光滑结构，然后构造局部平凡化并证明它们具有所需性质。下面的引理提供了一条捷径：它表明只要局部平凡化彼此之间以光滑转移函数重叠，构造局部平凡化就足够了。（此结果的更强形式另见习题~10-6。）

\begin{lemma}[向量丛坐标图引理]\label{lemma:ch10-向量丛坐标图引理}
设 $M$ 为带边或不带边的光滑流形，假设对每个 $p\in M$ 给定一个固定维数 $k$ 的实向量空间 $E_p$。令 $E=\coprod_{p\in M}E_p$，并令 $\pi\colon E\to M$ 为把 $E_p$ 中每个元素映到点 $p$ 的映射。此外假设我们给定下列数据：
\begin{enumerate}[(i)]
\item $M$ 的一个开覆盖 $\{U_\alpha\}_{\alpha\in A}$；
\item 对每个 $\alpha\in A$，一个双射 $\Phi_\alpha\colon\pi^{-1}(U_\alpha)\to U_\alpha\times\mathbb{R}^k$，它在每个 $E_p$ 上的限制是从 $E_p$ 到 $\{p\}\times\mathbb{R}^k\cong\mathbb{R}^k$ 的向量空间同构；
\item 对每对满足 $U_\alpha\cap U_\beta\neq\varnothing$ 的 $\alpha,\beta\in A$，一个光滑映射 $\gamma_{\alpha\beta}\colon U_\alpha\cap U_\beta\to GL(k,\mathbb{R})$，使得从 $(U_\alpha\cap U_\beta)\times\mathbb{R}^k$ 到其自身的映射 $\Phi_\alpha\circ\Phi_\beta^{-1}$ 具有形式
\begin{align}
\Phi_\alpha\circ\Phi_\beta^{-1}(p,v)=(p,\gamma_{\alpha\beta}(p)v).\label{align:ch10-转移协调}
\end{align}
\end{enumerate}
则 $E$ 具有唯一的拓扑与光滑结构，使其成为带边或不带边的光滑流形，并且是 $M$ 上秩为 $k$ 的光滑向量丛，以 $\pi$ 为投影、以 $\{(U_\alpha,\Phi_\alpha)\}$ 为光滑局部平凡化。
\end{lemma}

\begin{proof}
对每个点 $p\in M$，选取某个包含 $p$ 的 $U_\alpha$；选取 $M$ 的一个光滑坐标图 $(V_p,\varphi_p)$ 使得 $p\in V_p\subseteq U_\alpha$；令 $\widehat V_p=\varphi_p(V_p)\subseteq\mathbb{R}^n$ 或 $\mathbb{H}^n$（其中 $n$ 为 $M$ 的维数）。定义映射 $\widetilde\varphi_p\colon\pi^{-1}(V_p)\to\widehat V_p\times\mathbb{R}^k$ 为 $\widetilde\varphi_p=(\varphi_p\times\mathrm{Id}_{\mathbb{R}^k})\circ\Phi_\alpha$：
\begin{align*}
\pi^{-1}(V_p)\xrightarrow{\Phi_\alpha}V_p\times\mathbb{R}^k\xrightarrow{\varphi_p\times\mathrm{Id}_{\mathbb{R}^k}}\widehat V_p\times\mathbb{R}^k.
\end{align*}
我们将证明所有这种坐标图
\begin{align*}
\{(\pi^{-1}(V_p),\widetilde\varphi_p):p\in M\}
\end{align*}
满足光滑流形坐标图引理（\cref{lemma:ch01-35}）或其带边流形对应情形（习题~1.43）的条件，从而赋予 $E$ 一个带边或不带边光滑流形的结构。

作为双射映射的复合，$\widetilde\varphi_p$ 是到 $\mathbb{R}^n\times\mathbb{R}^k=\mathbb{R}^{n+k}$ 或 $\mathbb{H}^n\times\mathbb{R}^k\subseteq\mathbb{H}^{n+k}$ 中开子集上的双射。对任意 $p,q\in M$，容易验证
\begin{align*}
\widetilde\varphi_p\big(\pi^{-1}(V_p)\cap\pi^{-1}(V_q)\big)=\varphi_p(V_p\cap V_q)\times\mathbb{R}^k,
\end{align*}
这是开集，因为 $\varphi_p$ 是到 $\mathbb{R}^n$ 或 $\mathbb{H}^n$ 中开子集上的同胚。在两个这种坐标图重叠之处，我们有
\begin{align*}
\widetilde\varphi_p\circ\widetilde\varphi_q^{-1}=(\varphi_p\times\mathrm{Id}_{\mathbb{R}^k})\circ\Phi_\alpha\circ\Phi_\beta^{-1}\circ(\varphi_q\times\mathrm{Id}_{\mathbb{R}^k})^{-1}.
\end{align*}
由于 $\varphi_p\times\mathrm{Id}_{\mathbb{R}^k}$、$\varphi_q\times\mathrm{Id}_{\mathbb{R}^k}$ 以及 $\Phi_\alpha\circ\Phi_\beta^{-1}$ 都是微分同胚，该复合是微分同胚。因此满足\cref{lemma:ch01-35}~的条件~(i)--(iii)。由于开覆盖 $\{V_p:p\in M\}$ 有可数子覆盖，条件~(iv)~亦满足。

为检验 Hausdorff 条件~(v)，只需注意同一空间 $E_p$ 中任意两点都落在我们构造的某个坐标图中；而若 $\tilde\omega\in E_p$ 与 $\hat\omega\in E_q$ 且 $p\neq q$，我们可以选取 $p$ 与 $q$ 的不交邻域 $V_p$ 与 $V_q$，于是集合 $\pi^{-1}(V_p)$ 与 $\pi^{-1}(V_q)$ 是分别包含 $\tilde\omega$ 与 $\hat\omega$ 的不交坐标邻域。这样我们就赋予了 $E$ 带边或不带边光滑流形的结构。

关于此结构，每个映射 $\Phi_\alpha$ 都是微分同胚，因为就 $E$ 的坐标图 $(\pi^{-1}(V_p),\widetilde\varphi_p)$ 与 $V_p\times\mathbb{R}^k$ 的坐标图 $(V_p\times\mathbb{R}^k,\varphi_p\times\mathrm{Id}_{\mathbb{R}^k})$ 而言，$\Phi_\alpha$ 的坐标表示是恒等映射。$\pi$ 关于 $E$ 的同一坐标图与 $M$ 的坐标图 $(V_p,\varphi_p)$ 的坐标表示是 $\pi(x,v)=x$，因此 $\pi$ 也是光滑的。由于每个 $\Phi_\alpha$ 把 $E_p$ 映到 $\{p\}\times\mathbb{R}^k$，立即有 $\pi_1\circ\Phi_\alpha=\pi$，且由假设 $\Phi_\alpha$ 在纤维上是线性的。于是 $\Phi_\alpha$ 满足光滑局部平凡化的全部条件。

这是唯一的这种光滑结构，这一事实容易由“$\Phi_\alpha$ 到其像上为微分同胚”这一要求推出：任何满足相同条件的光滑结构必须包含我们构造的所有坐标图，因而等于这一个。
\end{proof}

下面给出一些例子，说明如何用坐标图引理从旧向量丛构造新向量丛。

\begin{example}[Whitney 和]\label{example:ch10-whitney和}
给定光滑流形 $M$ 以及 $M$ 上秩分别为 $k'$ 与 $k''$ 的光滑向量丛 $E'\to M$ 与 $E''\to M$，我们将在 $M$ 上构造一个新的向量丛，称为 $E'$ 与 $E''$ 的\textbf{Whitney 和}，它在每个 $p\in M$ 处的纤维是直和 $E'_p\oplus E''_p$。全空间定义为 $E'\oplus E''=\coprod_{p\in M}(E'_p\oplus E''_p)$，配以显然的投影 $\pi\colon E'\oplus E''\to M$。对每个 $p\in M$，选取 $p$ 的一个足够小的邻域 $U$，使得 $E'$ 与 $E''$ 分别在 $U$ 上有局部平凡化 $(U,\Phi')$ 与 $(U,\Phi'')$，并定义 $\Phi\colon\pi^{-1}(U)\to U\times\mathbb{R}^{k'+k''}$ 为
\begin{align*}
\Phi(v',v'')=\big(\pi'(v'),\,(\pi_{\mathbb{R}^{k'}}\circ\Phi'(v'),\,\pi_{\mathbb{R}^{k''}}\circ\Phi''(v''))\big).
\end{align*}
假设我们另有一对局部平凡化 $(\widetilde U,\widetilde\Phi')$ 与 $(\widetilde U,\widetilde\Phi'')$。令 $\gamma'\colon(U\cap\widetilde U)\to GL(k',\mathbb{R})$ 与 $\gamma''\colon(U\cap\widetilde U)\to GL(k'',\mathbb{R})$ 为相应的转移函数。则 $E'\oplus E''$ 的转移函数具有形式
\begin{align*}
\widetilde\Phi\circ\Phi^{-1}\big(p,(v',v'')\big)=(p,\gamma(p)(v',v'')),
\end{align*}
其中 $\gamma(p)=\gamma'(p)\oplus\gamma''(p)\in GL(k'+k'',\mathbb{R})$ 是分块对角矩阵
\begin{align*}
\begin{pmatrix}\gamma'(p)&0\\ 0&\gamma''(p)\end{pmatrix}.
\end{align*}
由于它光滑地依赖于 $p$，由坐标图引理可知 $E'\oplus E''$ 是 $M$ 上的光滑向量丛。
\end{example}

\begin{example}[向量丛的限制]\label{example:ch10-丛的限制}
设 $\pi\colon E\to M$ 为秩 $k$ 向量丛，$S\subseteq M$ 为任意子集。我们把 $E$ 在 $S$ 上的\textbf{限制}定义为集合 $E|_S=\coprod_{p\in S}E_p$，配以通过限制 $\pi$ 得到的投影 $E|_S\to S$。若 $\Phi\colon\pi^{-1}(U)\to U\times\mathbb{R}^k$ 是 $E$ 在 $U\subseteq M$ 上的局部平凡化，则它限制为一个双射 $\Phi|_U\colon(\pi|_S)^{-1}(U\cap S)\to(U\cap S)\times\mathbb{R}^k$，容易验证这些构成 $E|_S$ 上向量丛结构的局部平凡化。若 $E$ 是光滑向量丛，而 $S\subseteq M$ 是浸入或嵌入子流形，则由坐标图引理易知 $E|_S$ 是光滑向量丛。特别地，若 $S\subseteq M$ 是光滑（嵌入或浸入）子流形，则限制丛 $TM|_S$ 称为 $M$ 上的\textbf{环境切丛}。
\end{example}
\end{document}
